\documentclass[conference,10pt]{IEEEtran}
\pdfmapfile{+ucr.map}
\usepackage{amsmath,amssymb,booktabs,graphicx}
\usepackage[hidelinks]{hyperref}
\hypersetup{pdfauthor={},pdftitle={TilePlace: Deterministic Transactional Placement Refinement with a Full-Color Hardware Kernel}}
\newcommand{\system}{\textsc{TilePlace}}
\newcommand{\vpr}{\textsc{VPR}}
\setlength{\emergencystretch}{1em}

\begin{document}
\title{TilePlace: Deterministic Transactional Placement Refinement with a Full-Color Hardware Kernel}
\author{Anonymous authors}
\maketitle

\begin{abstract}
TilePlace refines legal FPGA placements using deterministic transactions and
a native routed acceptance oracle. It separates integer candidate selection
from routing decisions, preserves complete macro identities, and implements
a color invocation as an explicit-array hardware kernel. In a source-sealed
13-design CPU cohort, nine designs complete paired evaluation; their geometric
mean critical-path delay and wirelength changes are $-1.944\%$ and $-0.213\%$.
Four failures remain in the denominator. The verified hardware baseline passes
nine RTL transactions and a nine-case native differential regression. Internal
4\,ns timing closure remains unresolved. These results establish the current
functional scope, without claiming physical-board speedup, power, or energy.
\end{abstract}

\begin{IEEEkeywords}
FPGA placement, deterministic transactions, routed feedback, high-level synthesis.
\end{IEEEkeywords}

\noindent\textbf{Working draft, not submission-ready.}
Final same-source timing closure, implementation adoption, and manuscript
acceptance remain pending. Numerical claims below distinguish the completed
CPU cohort, the verified hardware baseline, and diagnostic timing candidates.

\section{Introduction}

FPGA placement must respect heterogeneous sites, subtiles, fixed blocks,
regions, and macros while controlling routed delay and congestion.
VPR supplies a native placement, routing, and timing reference for this
problem~\cite{betz1997vpr,vtr8}. A local geometric improvement, however, does
not by itself demonstrate better routed timing. A proposal can reduce a net's
bounding box while lengthening its critical source--sink connection, or can
displace a macro member into an illegal site.

\system{} addresses the post-placement refinement problem. Starting from a
legal packed placement, it enumerates bounded local transactions, ranks them
with integer costs, and constructs a deterministic nonconflicting batch.
Native routing and timing analysis decide whether that batch becomes the
new incumbent. Rejection restores the pre-trial placement exactly; stable
bisection and complete rejected identities permit bounded exploration of
alternative proposals.

This separation also defines a hardware boundary. A color computation may
evaluate and select transactions while the host retains expensive routed
acceptance and feedback. The challenge is to preserve complete transaction
identities, shared-snapshot semantics, and conflict sets across an explicit
memory interface. Checking only final coordinates or a scalar cost can miss
errors in scratch records consumed by later stages.

The work contributes a deterministic transaction contract for routed-oracle
refinement; a full-color integer kernel and native adapter implementing that
contract; and an evidence protocol linking complete writable-port RTL checks,
native search trajectories, exact input identities, and physical reports to
specific source snapshots. The present evaluation supports functional
equivalence and a limited CPU QoR result. It does not yet establish internal
4\,ns closure or application acceleration on a physical FPGA.

\section{Deterministic Routed-Oracle Refinement}

\subsection{Snapshots, proposals, and objective}

A block location includes coordinates, layer, and subtile. Macros consist of
a root and members with fixed relative offsets. Each candidate materializes
all member destinations before checking type compatibility, occupancy,
fixed or regional constraints, and macro shape. A failed member check rejects
the complete candidate; members cannot commit independently.

The evaluated protocol partitions the device into $2\times2$ pages with four
colors. A phase freezes positions, occupancy, objective inputs, and feedback.
An initiating unit searches its page and neighboring pages with a bound of
64 target sites and stable coordinate and identifier tie breaks. Occupied
targets produce swaps and chains of two or three units; compatible macros
may swap as complete objects. The routed protocol focuses initiating units
on blocks in the 20 worst current setup paths. Helper units can still
participate in a transaction.

Page coloring establishes a deterministic schedule, not independence:
long nets and macros can cross page boundaries. For each transaction $T$,
the computation records the old/new occupied-site union $S(T)$ and all
touched nets $N(T)$. It recomputes objective changes against the same frozen
snapshot. The wire term uses touched-net bounding boxes. A connection-aware
timing term retains sink identity:
\begin{equation}
C(P)=\sum_{c=(s,t)}q_c\bigl(|x_s-x_t|+|y_s-y_t|\bigr),
\end{equation}
where $q_c$ is a quantized criticality in $[0,65535]$.
Routed setup-path and congestion feedback augment proposal ranking.

The source-sealed cohort's initial timing weights come from a fixed CPU path
generator and a patched timing exporter supplied with its quantized edge
delays. They are proxy-derived weights, not measurements of native baseline
route delay. The native routed oracle remains the final QoR criterion.

\subsection{Selection and legality}

Candidates are ordered by improvement and stable transaction identity.
Scanning that order accepts $T$ only if
\begin{equation}
S(T)\cap S_A=\emptyset\quad\text{and}\quad N(T)\cap N_A=\emptyset,
\end{equation}
where $S_A,N_A$ accumulate previously accepted transactions. The explicit
kernel also carries unit conflict ranges. Complete identities include
operation kind, nested member contents, and trial coordinates, rather than
relying on pool positions or host indices.

For a fixed input contract, enumeration and integer evaluation yield one
ordered candidate sequence. An induction over the conflict scan gives one
accepted batch. Because accepted transactions have disjoint old/new site
sets and touched nets, each retains the before-state used during evaluation,
and applying the batch is equivalent to its deterministic serial order.
Member legality checks preserve macro structure and occupancy. These
arguments assume valid input backing and are supplemented by native
placement and routing validation.

\subsection{Routed acceptance and bounded exploration}

Let $P_0$ be the baseline and $P_i$ the current incumbent. A proposed $P'$ is
accepted only if native placement, routing, and timing checks succeed,
\begin{equation}
WL(P')\le1.005\,WL(P_0),
\end{equation}
and either $CPD(P_i)-CPD(P')\ge0.001$\,ns or CPD is exactly equal and routed
wirelength strictly decreases. The cap is baseline-relative, preventing
successive timing improvements from accumulating unlimited wire cost.

Before a trial, the controller snapshots all positions. Rejection restores
them exactly and splits a non-atomic batch at its stable midpoint, trying
halves in fixed order while budget remains. A rejected atomic identity enters
an incumbent-specific forbidden set. Later visits consider the next-ranked
available proposal. Acceptance clears this set and refreshes native setup-path
and congestion feedback because the incumbent has changed.

The protocol allows four epochs and 24 candidate closures. Every failed
closure consumes budget. With fixed routing binary, seed, channel width,
and options, the resulting controller trace can be checked transaction by
transaction, including placement hashes, decisions, restoration, and feedback.
Determinism does not imply that every local optimum is explored or that a
bounded run finds the best placement.

% Working section: structural description of the verified October 4 baseline.
% Final adoption and timing results must be bound before submission.
\section{Full-Color Hardware Computation}
\label{sec:hardware}

\subsection{Boundary between local selection and routed acceptance}

The hardware computation implements one color invocation of the integer
transaction algorithm. Its inputs are an immutable placement snapshot,
unit and member metadata, objective data, configuration, and forbidden
transaction identities. Its outputs are selected transactions, updated local
positions, diagnostics, and per-kind counts. The host retains the epoch and
color schedule, native VPR routing and timing analysis, closure budget,
rollback, batch bisection, and feedback refresh. Consequently, selection by
the color computation is a proposal to the routed controller; it does not
establish an improvement in routed critical-path delay.

The native integration enforces this distinction by invoking the shared
integer implementation on a private copy of the position array. It decodes
the selected records into the proposal protocol consumed by the existing
controller. A rejected native closure therefore restores the incumbent at
the host boundary, independently of the local position updates used to
construct a proposal. The native integration is an implementation equivalence
path, not a measurement of host-to-device transfers or board execution.

\begin{figure}[t]
\centering
\small
\begin{tabular}{@{}c@{}}
Host snapshot and forbidden identities\\
$\downarrow$\\
\fbox{Eligibility and five-kind enumeration}\\
$\downarrow$\\
\fbox{Trial legality and integer evaluation}\\
$\downarrow$\\
\fbox{External filtering and identity encoding}\\
$\downarrow$\\
\fbox{Per-group best, deduplication, conflicts}\\
$\downarrow$\\
\fbox{Batch validation, local commit, statistics}\\
$\downarrow$\\
Host proposal decoding and native routed oracle
\end{tabular}
\caption{Functional order within one color invocation and its host boundary.
The boxes do not denote concurrent streaming stages or an achieved
initiation interval.}
\label{fig:hardware-color}
\end{figure}

\subsection{Explicit identities and conflict sets}

The top-level computation enumerates moves, swaps, macro swaps, and chains
of two or three units. Descriptors retain their indices throughout
evaluation, including unchanged and rejected candidates with zero score.
Canonical identities encode the operation kind, complete nested unit members,
and trial coordinates. External forbidden identities are compared across
their respective pools using complete identity contents. Pool offsets or
host unit numbers alone are insufficient to identify the same transaction.

For each descriptor, evaluation derives touched-net ranges and the sorted,
deduplicated union of old and trial occupied sites. Both sides of a move
belong to its conflict set: checking only destination sites would omit
dependencies created by vacating a location. Explicit offsets carry the
descriptor's ranges through evaluation and encoding. Selection applies the
stable ranking, per-group choice, identity deduplication, and conflict checks
before any local position commit.

The commit operation uses two passes. The first validates the entire selected
batch, including record ranges, known block identifiers, and repeated blocks.
Only after these checks succeed does the second pass write positions.
Evaluation itself leaves positions unchanged. This contract assumes canonical
input backing and pairwise disjoint scratch buffers; scratch from a failed
invocation is discarded. It provides algorithmic batch validation, without
claiming a fault-tolerant memory transaction mechanism.

\subsection{Word interface and capacity}

The synthesized top exposes explicit arrays with 8-, 32-, and 64-bit element
types, 51 logical AXI master bundles, and an AXI-Lite control interface.
Configuration occupies indices 0--35. Fourteen result words report stage and
status, enumeration and movement counts, and the five accepted-kind counts.
The complete C/RTL comparison checks all 30 writable ports, including scratch
arrays and positions, rather than checking only a final objective value.

Logical master bundles do not specify physical HBM channel assignments.
The current out-of-context (OOC) implementation does not establish platform
arbitration, bandwidth, host-device transfer cost, or application throughput.
Likewise, the fixture's \texttt{depth=4096} interface declarations are not proof
that arbitrary benchmark instances fit. The deployment contract must also
account for input counts, checked ranges, scratch capacities, host allocation,
and the eventual memory mapping.

\subsection{Preserving semantics while changing scheduling}

The timing experiments change scheduling and function sharing while retaining
integer widths and transaction semantics. An explicit C++ temporary is not
sufficient evidence of a cycle boundary: HLS may retain a combinational
same-state bypass. Each boundary experiment therefore inspects the generated
RTL, identifies the capture and consuming states, and binds that observation
to the exact generated source. A separate experiment inlines trial-location
decoding to reduce a shared address-selection path; its structural check
examines the resulting module and address logic instead of claiming a new
register boundary.

The physical flow reads a 4\,ns clock constraint before synthesis for
\texttt{xcu280-fsvh2892-2L-e}, then places, optimizes, and routes the generated
design. Acceptance requires the same candidate to pass full functional
regression and final setup, hold, pulse-width, internal-constraint, routing,
and DRC checks. The process must actually exit successfully, with frozen
source and output hashes verified. HLS frequency estimates and intermediate
placement estimates cannot substitute for the final routed report. Because
this OOC flow supplies no input/output delays, even a passing result would
establish internal timing only; physical U280 performance requires separate
platform integration and execution evidence.


\section{Evaluation Protocol}

\subsection{Separate questions and evidence}

We distinguish four questions. CPU QoR asks whether the routed controller
improves paired native placements. Native differential regression asks whether
the shared integer implementation preserves search and feedback behavior.
C/RTL co-simulation asks whether the explicit AXI top matches complete
expected output arrays for its fixtures. Physical implementation asks whether
that specific generated design satisfies internal timing and routing gates.
A passing result at one boundary cannot replace another.

Experiments use Ubuntu 22.04 and Intel Xeon Gold 6530 hosts. Native closures
use VTR revision \texttt{8cb20aa5293b5ce5}\allowbreak
\texttt{f6a05c88b37b38e1305d9b07}; binary hashes accompany the records.
Hardware generation and implementation use Vitis HLS and Vivado 2022.2.
The OOC target is XCU280 at 4\,ns, with the clock constraint read before
synthesis. No physical U280 execution is available in the recorded environment.

\subsection{Source-sealed CPU cohort}

Before refinement outcomes, the execution protocol seals 13 local designs,
source and binary identities, architecture, seed 1, channel width 200,
objective weights 2/2/1, four epochs, 24 closures, and 20 feedback paths.
Baseline and candidate use the same packed netlist and routing options with
one router worker. Each successful design has two complete repetitions.
There are no outcome-based parameter changes or retries. Geometric means
use candidate/baseline ratios over successful pairs; failures remain visible
in the registered denominator.

Native baseline runs generate placement, macros, and routing resources.
The input path generator has a fixed 30-iteration limit. Failure to produce
a valid input ends that case before refinement. This distinction prevents
input failures from being described as unsuccessful placement optimizations.
The local exposure review is incomplete, so source sealing does not establish
strict historical blindness or broad architecture generalization.

\subsection{Functional hardware baseline}

The verified baseline has two distinct RTL suites. Five transactions exercise
the five operation kinds; four additional transactions include external
forbidden fallback behavior. Each transaction compares all 30 writable ports,
giving 150 and 120 complete typed-array comparisons. Acceptance also requires
C postchecking, a co-simulation PASS, an actual zero process return code,
and unchanged frozen sources. These are fixture-level equivalence checks,
not exhaustive proofs for arbitrary instances or malformed memories.

The native regression executes nine successful cohort cases through the
integer adapter. Independent checks compare complete search and feedback
signatures and final placements, verify commands against the reference
configuration, and match all actual input identities to the references.
This regression contains 202 native trials and 88 checked color calls with
129 unique frozen input files. It is separate from the CPU cohort's 404
trials across two repetitions; adding these counts would overstate coverage.

\section{Results and Current Limits}

\subsection{Paired routed QoR}

\begin{table}[t]
\centering\small
\caption{All registered CPU cases. Geometric means use nine successful pairs;
four failures remain in the cohort denominator.}
\label{tab:cohort}
\begin{tabular}{lrrrr}
\toprule
Design & Base (ns) & Final (ns) & $\Delta$CPD & $\Delta$WL \\
\midrule
adder & \multicolumn{4}{c}{baseline failed} \\
alu4 & 5.96812 & 5.96812 & +0.00\% & -0.20\% \\
apex2 & 6.69869 & 6.43134 & -3.99\% & +0.17\% \\
des & 5.01530 & 5.01530 & +0.00\% & -0.45\% \\
ex1010 & \multicolumn{4}{c}{input failed} \\
ex5p & 5.88580 & 5.61423 & -4.61\% & +0.11\% \\
misex3 & 5.08031 & 5.03800 & -0.83\% & -0.52\% \\
s298 & \multicolumn{4}{c}{input failed} \\
s38417 & \multicolumn{4}{c}{input failed} \\
s38584.1 & 6.90921 & 6.76483 & -2.09\% & -0.21\% \\
seq & 5.45825 & 5.45715 & -0.02\% & -0.13\% \\
spla & 7.65176 & 7.64733 & -0.06\% & +0.10\% \\
tseng & 6.64179 & 6.26454 & -5.68\% & -0.78\% \\
\midrule
Geo. (9/13) & -- & -- & -1.94\% & -0.21\% \\
\bottomrule
\end{tabular}

\end{table}

Nine of 13 registered designs complete evaluation (Table~\ref{tab:cohort}).
The independent audit verifies both repeated trajectories and 404 native
trials. Seven designs improve CPD and two preserve it. The successful-subset
geometric-mean CPD and routed wirelength changes are $-1.944\%$ and
$-0.213\%$, respectively. The largest observed CPD changes are
$-5.680\%$ for \texttt{tseng}, $-4.614\%$ for \texttt{ex5p}, and
$-3.991\%$ for \texttt{apex2}. These are paired post-route CPU results, not
hardware throughput measurements.

\texttt{adder} lacks a required architecture model. \texttt{ex1010},
\texttt{s298}, and \texttt{s38417} have legal native baselines but retain
8, 2, and 2 overused resources in the CPU input path generator after its
30-iteration limit. Refinement is not launched for these cases. Their
exclusion from the geometric mean is an execution limitation, not a claim
that the method would improve them.

\subsection{Suppression ablation and provenance}

A separate four-design fixed-source ablation compares rejected-identity
suppression enabled and disabled under the same input and feedback parser.
It finds no CPD advantage. In \texttt{diffeq}, suppression changes the trace
from eight to sixteen closure calls and enables one accepted transaction,
reducing wirelength by $0.055\%$ relative to disabled suppression. The other
three cases preserve final QoR; paired geometric-mean wirelength changes by
$-0.014\%$. Thus the mechanism can expose an alternative transaction, but
this experiment does not establish a general timing benefit.

Earlier archived cohorts have a timing-parser provenance gap: a fixed-source
replay reconstructs a different initial feedback edge set from the same
timing report. Their recorded outcomes are retained in the artifact history,
but are not pooled into the source-sealed result above or used to support
a single-frozen-implementation claim.

\subsection{Functional correctness and internal timing}

The verified October 4 baseline passes both full RTL suites and the native
regression described above. With the 4\,ns constraint applied before
synthesis, its routed OOC implementation has WNS $-0.660$\,ns, TNS
$-1455.593$\,ns, and 10,507 failing setup endpoints. All 335,892 routable
nets are routed and hold checks pass; setup timing does not.

An independent experiment starts from that baseline's archived routed
checkpoint and applies post-route physical optimization followed by routing
without timing relaxation. WNS improves to $-0.589$\,ns, but TNS remains
$-1450.798$\,ns with 10,506 failing setup endpoints. All 335,894 nets are
routed and hold checks pass. This experiment preserves the baseline's
logical implementation and original checkpoint; it still does not meet the
4\,ns gate.

\begin{table}[t]
\centering\small
\setlength{\tabcolsep}{3pt}
\caption{Completed October 12--14 internal 4\,ns OOC experiments. Each row
is a separate candidate and all fail setup timing. WNS and TNS are in ns;
fails counts setup endpoints. Rows are diagnostic
experiments, not a controlled causal ablation.}
\label{tab:ooc}
\begin{tabular}{lrrr}
\toprule
Candidate & WNS & TNS & Fails \\
\midrule
Owners return wait & -0.584 & -1188.139 & 8326 \\
Trial equality protocol & -0.480 & -510.018 & 4384 \\
Commit find comparison & -0.676 & -3395.322 & 16618 \\
Trial offset validation & -0.563 & -993.513 & 7455 \\
Best unit validation & -0.422 & -404.283 & 3630 \\
Identity kind count & -0.417 & -356.353 & 3702 \\
Translation index return & -0.707 & -1942.742 & 11072 \\
Identity kind nonzero & -0.627 & -810.712 & 5748 \\
Roots shift end & -0.237 & -59.007 & 918 \\
Eligible validation index & -0.509 & -784.205 & 5140 \\
Trial location inline & -0.632 & -1026.927 & 6444 \\
Encode status return & -0.229 & -74.812 & 1229 \\
Best dimension guard & -0.153 & -10.711 & 328 \\
Roots page guard & -0.603 & -572.762 & 4642 \\
Identity member nonzero & -1.415 & -3398.806 & 14019 \\
Cost position compare & -0.267 & -180.754 & 2496 \\
Selection member return & -0.277 & -104.519 & 1243 \\
Commit find return & -0.464 & -684.959 & 6226 \\
Position entry ID & -0.346 & -170.455 & 2067 \\
Position condition return & -0.479 & -420.669 & 4513 \\
Position remaining span & -0.388 & -88.520 & 1162 \\
Position ID compare return & -0.323 & -95.798 & 1380 \\
\bottomrule
\end{tabular}

\end{table}

Table~\ref{tab:ooc} retains 22 completed timing candidates. The best WNS in
this set is $-0.153$\,ns for the best-dimension-guard candidate, with TNS
$-10.711$\,ns and 328 failing setup endpoints. Its 332,910 nets are fully
routed; hold and pulse-width checks pass, and DRC reports 25 warnings with
no critical warnings or errors. It still fails the timing gate. Moreover,
the baseline's correctness results cannot be transferred to this changed
candidate. A final implementation must satisfy both functional and physical
gates with its own source bindings.

A separate post-route optimization of the roots shift-end candidate improves
its WNS from $-0.237$ to $-0.170$\,ns, with TNS $-53.134$\,ns and 896 setup
failures. This physical variant also fails the unchanged 4\,ns gate.
Post-route optimization of best-dimension-guard reduces its WNS to
$-0.086$\,ns, with TNS $-7.674$\,ns and 292 setup failures. All 332,917
nets are routed and hold checks pass, but setup closure remains unresolved.
A subsequent AggressiveExplore pass on that checkpoint reaches WNS
$-0.073$\,ns and TNS $-7.164$\,ns, with 287 setup failures. Its 332,921
nets are fully routed, hold and pulse-width checks pass, and DRC reports no
critical warnings or errors. This improvement still does not meet 4\,ns.
A further fanout-oriented physical pass reaches WNS $-0.067$\,ns and TNS
$-5.634$\,ns, with 267 setup failures and 332,923 fully routed nets. Hold,
pulse-width, and DRC error checks pass; the setup timing gate remains open.
Alternate critical-cell replication from that checkpoint further improves
WNS to $-0.050$\,ns and TNS to $-4.143$\,ns, with 251 setup failures.
All 332,923 nets remain routed, and hold, pulse-width and DRC error checks
pass. An iterative AggressiveExplore pass then reaches WNS $-0.047$\,ns
and TNS $-3.488$\,ns, with 224 setup failures and 332,925 fully routed
nets. Hold, pulse-width and DRC error checks still pass, but setup timing
remains below the unchanged 4\,ns requirement for this logical candidate.

Intermediate HLS estimates and partially completed simulations do not close
that gap. Register-boundary inspection can reject an ineffective scheduling
change before implementation, but does not establish timing improvement.
For example, capturing the position ID before binary search produces the
intended RTL boundary, yet its routed WNS is $-0.346$\,ns, worse than its
best-dimension-guard parent's $-0.153$\,ns. Its worst reported path crosses
selection member-range checking and AXI read control. A separate registered
64-bit search-condition return also forms its intended RTL boundary, yet
reaches WNS $-0.479$\,ns with 4,513 setup failures. Its worst path crosses
enumeration control and the position-memory AXI read-buffer enable;
eight of the 20 reported worst paths still end at position-query register
enables. Replacing the position search interval with a full-width start and
remaining length removes the low/high loop comparison in generated RTL,
but reaches WNS $-0.388$\,ns with 1,162 setup failures. Its worst reported
path crosses identity control and the unit-members AXI read-buffer enable.
Registering the exact 32-bit position-ID predicate also forms its intended
boundary, but reaches WNS $-0.323$\,ns with 1,380 setup failures. Its worst
path crosses identity count comparison and commit return control; routing
accounts for 85.8\% of that path's data delay.
Local boundary formation therefore does not establish whole-kernel timing
closure.
No result here measures board latency, throughput, transfer costs, power,
or energy. A clock target alone cannot be converted into application speedup.

\subsection{Baseline resources and simulated fixture cost}

The verified baseline's routed report uses 180,887 CLB LUTs, 223,917 CLB
registers, 68.5 block-RAM tiles, 46 DSPs, and no UltraRAM. LUT and register
utilization are $13.88\%$ and $8.59\%$ of the target, respectively. These
counts belong to the baseline implementation above; they cannot be combined
with a different candidate's timing to describe one realized design.

\begin{table}[t]
\centering\small
\caption{Verified baseline fixture latency in XSim clock cycles. Each row
summarizes its entire suite, including the generated AXI memory model.}
\label{tab:cycles}
\begin{tabular}{lrr}
\toprule
Fixture suite & Min cycles & Max cycles \\
\midrule
Five-kind (5) & 3,281,728 & 14,121,019 \\
Fallback (4) & 654,891 & 720,319 \\
\bottomrule
\end{tabular}

\end{table}

Table~\ref{tab:cycles} reports the range of per-invocation latency from
completed co-simulation. Transaction records reproduce the summary minima,
maxima, integer averages, and total execution cycles. These fixtures exercise
complete color invocations, not isolated move-operation costs. Their ranges
do not bound benchmark workloads or later candidates. In particular, the
simulation's 4\,ns clock setting is not achieved physical timing; multiplying
cycles by that target would not establish board runtime. The large cycle
counts also preclude assuming high throughput from resource fit alone.

\section{Related Work}

VPR and VTR provide architecture-aware placement and routed evaluation
infrastructure~\cite{betz1997vpr,vtr8}. Analytical approaches address broader
placement and packing objectives: routing-architecture-aware
placement~\cite{chen2015routingaware}, UTPlaceF~\cite{li2018utplacef},
elfPlace~\cite{li2019elfplace}, DREAMPlaceFPGA~\cite{rajarathnam2022dreamplacefpga},
and AMF-Placer 2.0~\cite{liang2024amf}. FPGA-accelerated placement has also
been studied~\cite{dhar2019accelerated}. \system{} operates after a legal
packed placement and makes native routing an explicit acceptance boundary.
The present evidence does not provide a runtime or QoR comparison against
those distinct end-to-end flows. Its hardware contribution is the explicit
full-color transaction implementation and its verification boundary.

\section{Limitations and Reproducibility}

The successful CPU subset is small and logic-oriented, uses one routing seed
and fixed budgets, and depends on a proxy input generator with capacity
failures. Surrogate weights and local proposal bounds may exclude useful
transactions. Stable bisection and suppression improve exploration discipline
without guaranteeing optimality. A complete runtime study must separately
measure proposal computation, memory transfers, native routing, feedback, and
rollback; the current functional evidence cannot predict their balance.

The artifact separates frozen source, fixture definitions, native reference
inputs, commands, terminal return codes, full typed outputs, search traces,
and implementation reports. Table generation checks their recorded identities
and stops on mismatch. Failed timing variants and failed CPU cases remain
visible. Final release requires a single adopted source manifest linking the
passing RTL, native, and routed results, followed by a manuscript update and
complete anonymity and format checks. Those release conditions remain open.

\section{Conclusion}

\system{} separates deterministic local transaction selection from native
routed acceptance and implements the color computation across a fully
checked explicit-word interface. The source-sealed CPU cohort improves
successful-subset routed CPD by $1.944\%$, while the verified hardware
baseline passes the stated functional regressions. Internal 4\,ns timing
closure and final implementation adoption remain unresolved; physical-board
acceleration is not established.

\clearpage
\bibliographystyle{IEEEtran}
\bibliography{references}
\end{document}
